\documentclass[10pt, twocolumn, letterpaper]{article}
% =============================================
% PAQUETES
% =============================================
\usepackage[utf8]{inputenc}
\usepackage[T1]{fontenc}
\usepackage{mathptmx}        % Times New Roman (mejor soporte que el paquete 'times')
\usepackage{textcomp}
\usepackage{babel}
\usepackage{geometry}
\usepackage{titlesec}
\usepackage{booktabs}
\usepackage{tabularx}
\usepackage{array}
\usepackage{graphicx}
\usepackage{caption}
\usepackage{multicol}
\usepackage{amsmath}
\usepackage{hyperref}
\usepackage{xcolor}
\usepackage{enumitem}
\usepackage{float}
\usepackage{rotating}

% =============================================
% GEOMETRÍA DE PÁGINA (Formato A1 COMEA)
% =============================================
\geometry{
letterpaper,
top=19mm,
bottom=25.4mm,
left=17.3mm,
right=17.3mm,
columnsep=4.22mm
}
\pagestyle{empty}

% =============================================
% FORMATO DE SECCIONES
% =============================================
\titleformat{\section}
{\fontsize{10}{12}\selectfont\bfseries\centering}
{\Roman{section}.}
{0.5em}
{\MakeUppercase}

\titleformat{\subsection}
{\fontsize{10}{12}\selectfont\itshape}
{\Alph{subsection}.}
{0.5em}
{}

\titlespacing*{\section}{0pt}{6pt}{3pt}
\titlespacing*{\subsection}{0pt}{4pt}{2pt}

\renewcommand{\thesection}{\Roman{section}}
\renewcommand{\thesubsection}{\Alph{subsection}}

% =============================================
% CAPTIONS
% =============================================
\captionsetup{font=footnotesize, labelsep=space, justification=centering}
\captionsetup[table]{position=top, name=TABLA, labelsep=newline}
\captionsetup[figure]{position=bottom, name=Fig.}

\hypersetup{colorlinks=false}

% =============================================
% INICIO DEL DOCUMENTO
% =============================================
\begin{document}

% =============================================
% ENCABEZADO
% =============================================
\twocolumn[{%
\begin{center}
{\fontsize{9}{11}\selectfont Categoría: Ruta Sincronizada — COMEA 2026}
\vspace{6pt}

{\fontsize{24}{28}\selectfont\bfseries
Diseño, Construcción e Integración de un Sistema\\
de Drones para Ruta Sincronizada — COMEA 2026
}
\vspace{8pt}

{\fontsize{11}{13}\selectfont
Cabral Ávila D., Cervantes Urquidi O., Galindo Reyes M.~M.,
Soto Ramírez D.~E., Corral Gómez N.~I.,\\
Higuera Lugo B.~Y., Rivera Rubio B.~K., Rodríguez de Jesús A.~O., Rojas Puente J.~R.
}
\vspace{4pt}

{\fontsize{9}{11}\selectfont
Ingeniería Aeronáutica — Universidad Politécnica de Chihuahua (UPCH), México
}
\vspace{4pt}

{\fontsize{9}{11}\selectfont
Asesor técnico: Hidalgo Morales A. | Asesor académico: González Torres E.\\
Universidad Politécnica de Chihuahua, México
}
\vspace{10pt}
\rule{\linewidth}{0.4pt}
\vspace{6pt}
\end{center}
}]

% =============================================
% RESUMEN
% =============================================
\noindent{\fontsize{9}{11}\bfseries\itshape Resumen}—{\fontsize{9}{11}\itshape\selectfont
Este documento presenta el reporte técnico del equipo Black Bears de la
Universidad Politécnica de Chihuahua (UPCH) para la categoría de Ruta
Sincronizada de la competencia COMEA 2026. El proyecto contempla el diseño,
fabricación, integración y operación de dos vehículos aéreos no tripulados (UAV)
tipo cuadricóptero clase 5\,pulgadas, capaces de transportar de manera coordinada
un cubo de carga útil de $10\times10\times10$\,cm a través de un área de
$20\times40$\,m con obstáculos cilíndricos. La plataforma emplea controladores de
vuelo SpeedyBee F405~V3 con firmware Betaflight, motores EMAX ECO~II 2205~2300KV
y un mecanismo de sujeción compartido fabricado en PETG con varillas de fibra de
carbono. El proyecto es desarrollado por nueve estudiantes universitarios
multidisciplinarios durante veinte semanas, comprendidas entre mayo y octubre
de 2026.}
\vspace{4pt}

\noindent{\fontsize{9}{11}\bfseries\itshape Palabras clave:}
{\fontsize{9}{11}\itshape UAV, cuadricóptero, ruta sincronizada, Betaflight, mecanismo
de sujeción, coordinación multi-dron, COMEA 2026.}
\vspace{6pt}

% =============================================
% I. INTRODUCCIÓN
% =============================================
\section{Introducción}
La industria aeroespacial ha experimentado un crecimiento significativo en el uso
de sistemas aéreos no tripulados (UAV) para aplicaciones civiles, industriales y
científicas. La competencia COMEA 2026 representa una oportunidad para que
estudiantes de ingeniería desarrollen habilidades prácticas en diseño, manufactura,
integración y operación de plataformas aéreas.

El presente proyecto tiene como objetivo el diseño, construcción e integración de
dos UAV tipo cuadricóptero para participar en la categoría de Ruta Sincronizada,
en la que dos o más drones deben transportar de forma coordinada y teleoperada
un cubo de carga útil a través de un área de competencia con obstáculos. El
desarrollo será realizado por el equipo \textit{Black Bears} de la Universidad
Politécnica de Chihuahua, conformado por nueve estudiantes multidisciplinarios,
durante veinte semanas comprendidas entre el 25 de mayo y el 9 de octubre de 2026.

% =============================================
% II. INTEGRANTES
% =============================================
\section{Nombre y Funciones de Cada Integrante}
\subsection{Organización del Equipo}
El equipo \textit{Black Bears} está conformado por nueve integrantes con
responsabilidades claramente definidas, bajo la supervisión del asesor técnico
Andrés Hidalgo Morales y el asesor académico Edgar González Torres.

\begin{table}[H]
\caption{Integrantes y Roles del Equipo}
\fontsize{8}{9}\selectfont
\begin{tabularx}{\columnwidth}{>{\raggedright}p{2.4cm} X}
\toprule
\textbf{Integrante} & \textbf{Rol y Responsabilidades Técnicas} \\
\midrule
Daniela Cabral Ávila & Líder del proyecto. Coordinación general, Secciones I, II, XI. \\
\addlinespace
Oswaldo Cervantes Urquidi & Líder técnico y Análisis Estructural. Sección IV.F (FEA), Operaciones. \\
\addlinespace
Mia M. Galindo Reyes & Coordinadora de Electrónica y Control. Secciones VI.B, VI.D (PID, Algoritmos). \\
\addlinespace
Daniel E. Soto Ramírez & Responsable de Software y Piloto UAV 1. Secciones VI.B, VI.D. \\
\addlinespace
Noel I. Corral Gómez & Responsable de Propulsión y Piloto UAV 2. Sección V.C (Tracción). \\
\addlinespace
Brizeth Y. Higuera Lugo & Coordinadora de Diseño Mecánico. Secciones III.C, IV.A, IV.C, IV.E, VI.D. \\
\addlinespace
Brian K. Rivera Rubio & Responsable de Manufactura y Ensamble. Sección VII. \\
\addlinespace
Abdel O. Rodríguez de Jesús & Piloto Principal y Protocolos de Seguridad. Sección X. \\
\addlinespace
Juan R. Rojas Puente & Responsable de Documentación y Registro. Secciones VII.D, VIII.E. \\
\bottomrule
\end{tabularx}
\end{table}

\subsection{Cronograma de Trabajo}
El proyecto se desarrolla en ocho etapas durante veinte semanas. 

\begin{table}[H]
\caption{Etapas del Proyecto}
\fontsize{8}{9}\selectfont
\begin{tabularx}{\columnwidth}{>{\centering}p{0.5cm} >{\raggedright}p{2.2cm} X}
\toprule
\textbf{No.} & \textbf{Periodo} & \textbf{Etapa} \\
\midrule
1 & 25 may -- 12 jun & Organización y Planeación \\
2 & 15 -- 26 jun & Diseño Conceptual y Reporte Técnico \\
3 & 29 jun -- 10 jul & Adquisición de Componentes \\
4 & 13 -- 31 jul & Construcción e Integración \\
5 & 3 -- 21 ago & Configuración y Validación Inicial \\
6 & 24 ago -- 18 sep & Pruebas de Vuelo y Optimización \\
7 & 21 sep -- 2 oct & Entrenamiento y Simulación \\
8 & 5 -- 9 oct & Preparación Final y Competencia \\
\bottomrule
\end{tabularx}
\end{table}

% =============================================
% III. CONCEPTO DE DISEÑO
% =============================================
\section{Concepto de Diseño del Drone}
\subsection{Tipo de Drone}
Se emplea un cuadricóptero de configuración en~X, clase 5\,pulgadas (diagonal
máxima $\leq 40$\,cm), con cuatro rotores de paso fijo y control por variación
diferencial de velocidad. El peso estimado por unidad es de $\approx$\,350\,g
(sin carga útil), cumpliendo el límite de la categoría. Se utilizarán dos
unidades idénticas para la misión de Ruta Sincronizada.

\subsection{Arquitectura General del Sistema}
Cada UAV integra cuatro subsistemas principales:
\begin{enumerate}[leftmargin=*, itemsep=0pt, topsep=2pt]
\item \textbf{Propulsión:} 4 motores EMAX ECO~II 2205~2300KV + ESC 35A BLHeli\_32 4-en-1.
\item \textbf{Control:} FC SpeedyBee F405~V3 (Betaflight) + GPS uBlox M8N.
\item \textbf{Comunicación:} Radio RadiMaster TX16S / FlySky FS-i6X $\leftrightarrow$ receptor FS-iA6B / R81.
\item \textbf{Energía:} Batería LiPo 4S 1100\,mAh 100C (CNHL o Tattu).
\end{enumerate}
Adicionalmente, ambos drones comparten un mecanismo de sujeción de carga útil
(cubo $10\times10\times10$\,cm) fabricado en PETG con varillas de fibra de carbono.

\subsection{Diagrama de Bloques del Drone}
\noindent\fbox{\parbox{\dimexpr\columnwidth-2\fboxsep\relax}{%
\centering\fontsize{8}{10}\selectfont\itshape
[Diagrama de bloques pendiente — Entrega: 20 junio 2026]\\
Responsable: Noel Corral / Brizeth Higuera
}}

\subsection{Restricciones de la Competencia}
La Tabla~III resume los parámetros técnicos obligatorios de la categoría Ruta
Sincronizada según las bases de COMEA 2026.

\begin{table}[H]
\caption{Restricciones Técnicas — Ruta Sincronizada}
\label{tab:restricciones}
\fontsize{8}{9}\selectfont
\begin{tabularx}{\columnwidth}{>{\raggedright}p{2.8cm} X}
\toprule
\textbf{Parámetro} & \textbf{Especificación} \\
\midrule
Dimensión máxima & $\leq 40$\,cm en diagonal (clase 5\,") \\
Número de drones & 2--4 por equipo \\
Carga útil & Cubo de papel $10\times10\times10$\,cm, sujeción compartida \\
Altura de vuelo & 2.0--2.5\,m \\
Duración de misión & 5\,min continuos \\
Modo de operación & Teleoperado (pilotos humanos) \\
Área de operación & $20\times40$\,m con 4 obstáculos cilíndricos \\
\bottomrule
\end{tabularx}
\end{table}

\subsection{Justificación del Diseño Elegido}
Se seleccionó la clase 5\,pulgadas por ser la de mayor eficiencia
energética dentro del límite de 40\,cm, lo que permite alcanzar una relación
empuje-peso $\geq 3$:1 con componentes comerciales ampliamente disponibles en
México. La configuración de dos drones (mínimo reglamentario) reduce la
complejidad de coordinación respecto a tres o cuatro unidades.

El firmware Betaflight fue seleccionado por su madurez, comunidad de soporte y
capacidad de ajuste fino de parámetros PID, esenciales para la estabilidad bajo
carga dinámica compartida.

% =============================================
% IV. DISEÑO MECÁNICO
% =============================================
\section{Diseño Mecánico y Estructural}
\subsection{Diseño de la Estructura o Fuselaje}
\noindent\fbox{\parbox{\dimexpr\columnwidth-2\fboxsep\relax}{%
\centering\fontsize{8}{10}\selectfont\itshape
[En desarrollo — Diseño CAD del frame en fibra de carbono 5"]\\
Responsable: Brizeth Higuera | Fecha estimada: 10 julio 2026
}}

\subsection{Materiales Utilizados}
El frame principal se fabricará en fibra de carbono (laminado preimpregnado
2×2 twill), seleccionada por su alta relación resistencia/peso. El mecanismo
de sujeción de la carga útil se fabricará en PETG mediante impresión 3D, con
varillas de refuerzo en fibra de carbono de 3\,mm de diámetro.

\subsection{Distribución de Masas}
\noindent\fbox{\parbox{\dimexpr\columnwidth-2\fboxsep\relax}{%
\centering\fontsize{8}{10}\selectfont\itshape
[Cálculo de centro de masa pendiente — CAD en desarrollo]\\
Responsable: Brizeth Higuera
}}

\subsection{Tipo de Ensamble}
La unión entre brazos y placa central se realizará mediante tornillería M3 de
acero inoxidable. El mecanismo de sujeción se acopla al frame mediante
mosquetones de aluminio y cables Dyneema de 1\,mm (baja elasticidad).

\subsection{Planos e Isométrico en CAD}
\noindent\fbox{\parbox{\dimexpr\columnwidth-2\fboxsep\relax}{%
\centering\fontsize{8}{10}\selectfont\itshape
[Planos CAD (sistema ISO) pendientes]\\
Responsable: Brizeth Higuera | Software: SolidWorks/Fusion 360
}}

\subsection{Simulación Numérica de Esfuerzos}
\noindent\fbox{\parbox{\dimexpr\columnwidth-2\fboxsep\relax}{%
\centering\fontsize{8}{10}\selectfont\itshape
[Simulación FEA pendiente (ANSYS/CATIA)]\\
Responsable: Oswaldo Cervantes
}}

% =============================================
% V. PROPULSIÓN
% =============================================
\section{Sistema de Propulsión y Energía}
\subsection{Motores, Hélices y Controlador Electrónico de Velocidad}
Se utilizarán motores \textbf{EMAX ECO~II 2205 2300KV}, diseñados para
aplicaciones de racing y freestyle en clase 5\,". Sus características principales
son: resistencia interna $R_i \approx 0.08\,\Omega$, empuje máximo de
$\approx$\,750\,g por motor con hélice 5\,×\,4.3 en 4S, y corriente máxima
de 25\,A.

Las \textbf{hélices Gemfan 5143} (5\,×\,4.3, tripaleta) fueron seleccionadas
por su perfil de alto empuje a media velocidad, óptimo para vuelo de carga.

El \textbf{ESC 35A 4-en-1 BLHeli\_32} centraliza el control de los cuatro
motores en una sola placa, reduciendo masa y cableado.

\subsection{Baterías}
Se emplearán baterías \textbf{LiPo 4S 1100\,mAh 100C} (CNHL o Tattu).
Características principales: voltaje nominal 14.8\,V, corriente máxima de descarga: $C_{max} = 100 \times 1.1 = 110$\,A. Peso aproximado: 105\,g.

\subsection{Estimación de Tracción y Consumo}
El cálculo analítico completo se desarrollará una vez confirmados los datos de las fichas técnicas. Metodología a seguir:

\textbf{Empuje total estimado:}
\begin{equation}
T_{total} = 4 \times T_{motor}(RPM, \phi_{helice})
\end{equation}

\textbf{Relación empuje-peso (TWR):}
\begin{equation}
TWR = \frac{T_{total}}{(m_{dron} + m_{carga}) \cdot g} \geq 2.0
\end{equation}

\textbf{Autonomía estimada en hover:}
\begin{equation}
t_{hover} = \frac{C_{bat} \times \eta}{I_{hover}} \geq 5\,\text{min}
\end{equation}

\noindent\fbox{\parbox{\dimexpr\columnwidth-2\fboxsep\relax}{%
\centering\fontsize{8}{10}\selectfont\itshape
[Tabla de resultados numéricos y curva empuje vs. consumo: pendiente]\\
Responsable: Noel Corral | Fecha: 27 junio 2026
}}

\subsection{Justificación de Selección de Componentes}
La selección de motores 2205~2300KV responde a la necesidad de generar
suficiente empuje bajo carga compartida (cubo + mecanismo $\approx$\,200\,g adicionales por dron), manteniendo eficiencia energética para cumplir los 5\,min de misión.

% =============================================
% VI. SOFTWARE
% =============================================
\section{Software, Control y Navegación}
\subsection{Firmware Utilizado}
Se empleará \textbf{Betaflight} en el FC SpeedyBee F405~V3. Betaflight ofrece menor latencia en el lazo de control y soporte para GPS Rescue.

\subsection{Algoritmos de Control o Navegación}
\noindent\fbox{\parbox{\dimexpr\columnwidth-2\fboxsep\relax}{%
\centering\fontsize{8}{10}\selectfont\itshape
[Descripción de algoritmos de coordinación: en desarrollo]\\
Responsable: Daniel Soto / Mia Galindo
}}

\subsection{Modos de Vuelo}
Se utilizarán los modos \textit{Angle} (estabilizado) para la fase de carga y \textit{Horizon} para maniobras de baja velocidad. El modo GPS Rescue estará configurado como contingencia.

\subsection{Programación o Ajustes Realizados}
\noindent\fbox{\parbox{\dimexpr\columnwidth-2\fboxsep\relax}{%
\centering\fontsize{8}{10}\selectfont\itshape
[Parámetros PID y configuración de rates: pendiente]\\
Responsable: Daniel Soto / Mia Galindo
}}

\subsection{Diagramas de Flujo}
\noindent\fbox{\parbox{\dimexpr\columnwidth-2\fboxsep\relax}{%
\centering\fontsize{8}{10}\selectfont\itshape
[Diagramas de flujo de misión y procedimientos: en desarrollo]\\
Responsable: Brizeth Higuera
}}

\subsection{Justificación de Elementos Usados}
La elección de Betaflight frente a ArduPilot responde a la naturaleza teleoperada de la misión: Betaflight ofrece menor latencia en el lazo de control (típicamente $< 1$\,ms), favoreciendo la respuesta rápida ante correcciones del piloto.

% =============================================
% VII, VIII, IX
% =============================================
\section{Integración del Sistema}
\noindent\textit{Sección en desarrollo. Los contenidos de ensamblaje, problemas encontrados y material fotográfico serán documentados durante las Etapas 4 y 5 (13 julio -- 21 agosto 2026).}

\noindent\fbox{\parbox{\dimexpr\columnwidth-2\fboxsep\relax}{%
\centering\fontsize{8}{10}\selectfont\itshape
Responsable: Brian Rivera (ensamblaje) / Juan Rojas (fotografía)\\
Fecha estimada: 21 agosto 2026
}}

\section{Pruebas y Validación}
\noindent\textit{Sección en desarrollo. Las pruebas en tierra, de vuelo y sus resultados serán documentados durante la Etapa 6 (24 agosto -- 18 septiembre 2026).}

\noindent\fbox{\parbox{\dimexpr\columnwidth-2\fboxsep\relax}{%
\centering\fontsize{8}{10}\selectfont\itshape
Responsable: Oswaldo Cervantes (operaciones) / Juan Rojas (registro)\\
Fecha estimada: 18 septiembre 2026
}}

\section{Análisis de Fallas y Mejoras}
\noindent\textit{Sección en desarrollo. El análisis de fallas y limitaciones se completará tras las pruebas de vuelo.}

\noindent\fbox{\parbox{\dimexpr\columnwidth-2\fboxsep\relax}{%
\centering\fontsize{8}{10}\selectfont\itshape
Responsable: Oswaldo Cervantes / Daniel Soto\\
Fecha estimada: 25 septiembre 2026
}}

% =============================================
% X. SEGURIDAD
% =============================================
\section{Seguridad y Gestión de Riesgos}
\subsection{Medidas de Seguridad Implementadas}
\begin{itemize}[leftmargin=*, itemsep=0pt, topsep=2pt]
\item Uso de estuches anti-flama para almacenamiento y transporte de baterías LiPo.
\item Verificación de voltaje celda por celda antes de cada vuelo.
\item Zona de exclusión de 5\,m alrededor del área de vuelo durante pruebas.
\item Revisión de integridad estructural (hélices, tornillería) antes de cada sesión.
\end{itemize}

\subsection{Análisis de Riesgos}
\begin{table}[H]
\caption{Análisis de Riesgos}
\fontsize{8}{9}\selectfont
\begin{tabularx}{\columnwidth}{>{\raggedright}p{1.6cm} c c X}
\toprule
\textbf{Riesgo} & \textbf{P} & \textbf{I} & \textbf{Mitigación} \\
\midrule
Pérdida de señal radio & M & A & GPS Rescue activado; pilotos en LOS \\
Falla de batería en vuelo & B & A & Revisión pre-vuelo; baterías de respaldo \\
Colisión entre drones & M & A & Protocolos de vuelo sincronizado \\
Falla del mecanismo de sujeción & M & A & Pruebas de carga estática antes de vuelo \\
Daño estructural por crash & A & M & Frame y hélices de repuesto disponibles \\
\bottomrule
\end{tabularx}
\end{table}

\subsection{Procedimientos de Emergencia en Caso de Fallo}
\begin{enumerate}[leftmargin=*, itemsep=0pt, topsep=2pt]
\item \textbf{Pérdida de control:} Activar switch de desarm inmediatamente; el GPS Rescue tomará control si hay señal GPS.
\item \textbf{Incendio de batería:} Aislar con arena seca; no usar agua; evacuar el área.
\item \textbf{Falla del mecanismo de sujeción:} Descenso controlado coordinado entre ambos pilotos; posicionar la carga sobre zona segura.
\end{enumerate}

% =============================================
% XI. CONCLUSIONES
% =============================================
\section{Conclusiones}
\noindent\textit{Sección en desarrollo. Las conclusiones serán elaboradas tras la participación en COMEA 2026 (7 octubre 2026).}

\noindent\fbox{\parbox{\dimexpr\columnwidth-2\fboxsep\relax}{%
\centering\fontsize{8}{10}\selectfont\itshape
Responsable: Daniela Cabral\\
Fecha estimada: 9 octubre 2026
}}

% =============================================
% RECONOCIMIENTOS Y REFERENCIAS
% =============================================
\section*{Reconocimientos}
Los autores agradecen el apoyo de la Universidad Politécnica de Chihuahua (UPCH) por el respaldo institucional brindado al equipo \textit{Black Bears} para participar en COMEA 2026. Asimismo, se agradece la orientación del asesor técnico Andrés Hidalgo Morales y del asesor académico Edgar González Torres.

\begin{thebibliography}{9}
\fontsize{8}{9}\selectfont

\bibitem{comea2026}
COMEA, \textit{Bases y Reglamento de la Competencia de Drones COMEA 2026}, México, 2026.

\bibitem{betaflight}
Betaflight Development Team, \textit{Betaflight Configurator Documentation}, [Online]. Available: https://betaflight.com/docs

\bibitem{emax}
EMAX, \textit{ECO II Series Motor Datasheet — 2205 2300KV}, EMAX Technology, 2023. [Online]. Available: https://emaxmodel.com

\bibitem{speedybee}
SpeedyBee, \textit{F405 V3 Flight Controller User Manual}, SpeedyBee Technology, 2023. [Online]. Available: https://speedybee.com

\bibitem{ublox}
u-blox, \textit{NEO-M8 Series u-blox M8 Concurrent GNSS Modules Data Sheet}, u-blox AG, Thalwil, Switzerland, 2022.

\bibitem{mueller}
T. Mueller, \textit{Fundamentals of Unmanned Aircraft System Design}, AIAA Education Series, Reston, VA: American Institute of Aeronautics and Astronautics, 2019.

\bibitem{gemfan}
Gemfan, \textit{5143 Triblade Propeller Technical Specifications}, Gemfan Workshop, 2022.

\end{thebibliography}

\end{document}